% Figure: the lever rule for quality on a T-v (or p-v) two-phase tie line.
% A horizontal tie line at constant T spans from saturated liquid (v_f) to
% saturated vapour (v_g); a two-phase state of quality x sits at
% v = (1-x) v_f + x v_g, dividing the line so that the vapour fraction is
% the left segment over the whole. Schematic, with the dome for context.
\begin{figure}[htbp]
\centering
\begin{tikzpicture}
\begin{axis}[
  width=11cm, height=7.5cm,
  xlabel={specific volume $v$ (arb.\ units)},
  ylabel={temperature $T$ (arb.\ units)},
  xmin=0, xmax=10, ymin=0, ymax=7,
  axis lines=left,
  xtick=\empty, ytick=\empty,
  tick label style={font=\scriptsize},
  label style={font=\small},
  clip=true,
]
% dome
\addplot[engblue, very thick, smooth] coordinates {
  (1.4,1.0) (1.7,2.4) (2.3,3.8) (3.2,5.0) (4.3,5.6) (5.2,5.9)
};
\addplot[engblue, very thick, smooth] coordinates {
  (5.2,5.9) (6.2,5.6) (7.1,5.0) (7.9,4.0) (8.5,2.6) (8.9,1.0)
};
% tie line at T level y=3.4: from v_f on left limb to v_g on right limb
% left limb at y=3.4 -> x approx 2.15 ; right limb at y=3.4 -> x approx 8.25
\addplot[engred, very thick] coordinates {(2.15,3.4) (8.25,3.4)};
\addplot[only marks, mark=*, engblue, mark size=1.8pt] coordinates {(2.15,3.4) (8.25,3.4)};
% two-phase state at quality x=0.4: v = v_f + 0.4(v_g - v_f)
% v = 2.15 + 0.4*(8.25-2.15) = 2.15 + 2.44 = 4.59
\addplot[only marks, mark=square*, engorange, mark size=2.4pt] coordinates {(4.59,3.4)};
\node[engblue, font=\scriptsize, anchor=north] at (axis cs:2.15,3.3) {$v_f$};
\node[engblue, font=\scriptsize, anchor=north] at (axis cs:8.25,3.3) {$v_g$};
\node[engorange, font=\scriptsize, anchor=south] at (axis cs:4.59,3.5) {state, $x=0.4$};
% segment annotations: liquid-side length proportional to (1-x), vapour-side to x... actually
% the vapour mass fraction equals (v - v_f)/(v_g - v_f) = left segment / whole.
\draw[<->, enggray] (axis cs:2.15,2.9) -- (axis cs:4.59,2.9);
\node[enggray, font=\scriptsize, anchor=north] at (axis cs:3.37,2.85) {$\propto x$};
\draw[<->, enggray] (axis cs:4.59,2.9) -- (axis cs:8.25,2.9);
\node[enggray, font=\scriptsize, anchor=north] at (axis cs:6.42,2.85) {$\propto (1-x)$};
\end{axis}
\end{tikzpicture}
\caption[The lever rule for quality]{The lever rule for a two-phase
state, drawn on a temperature-volume diagram. At a fixed temperature
inside the dome the state lies on a horizontal tie line running from the
saturated-liquid volume $v_f$ to the saturated-vapour volume $v_g$. A
mixture of quality $x$ sits at $v=(1-x)v_f+xv_g$, so the vapour mass
fraction equals the length of the left segment (from $v_f$ to the state)
divided by the full tie-line length, and the liquid fraction equals the
right segment over the whole. The same construction reads quality from a
$p$-$v$ diagram and from the enthalpy, entropy, and internal-energy
weightings, because each two-phase property is the same linear average.
Coordinates are schematic.}
\label{fig:vol04ch02:quality-lever}
\end{figure}
